\section{Цветовые схемы: красный/чёрный, синий/белый, фиолетовый/золото, зелёный/чёрный, пастель}

Цветовая схема персонажа --- это не ``какие цвета красиво сочетаются''.
В аниме и character design цвет почти всегда работает как быстрая
психологическая грамматика. Он сообщает жанр, темперамент, опасность,
дистанцию, роль, возраст, уровень мистики, уровень милоты, уровень
``человечности'', а иногда даже моральную геометрию персонажа.

До того как зритель разберёт силуэт, волосы, глаза, одежду и голос, он
часто считывает цветовое пятно.

Красный/\hspace{0pt}чёрный --- опасность, страсть, агрессия, демоничность, власть,
запрет. Синий/\hspace{0pt}белый --- холод, чистота, рациональность, дистанция,
ангельскость, технология, больничность, winter/\hspace{0pt}AI. Фиолетовый/\hspace{0pt}золото ---
ночь, магия, интеллект, тайна, достоинство, внутреннее тепло.
Зелёный/\hspace{0pt}чёрный --- яд, ведьмовство, природа, вирус, хакерство,
токсичность, неестественность. Пастель --- moe, мягкость, доступность,
детскость, comfort, sweetness, idol/\hspace{0pt}mascot.

Это, конечно, не жёсткие правила. Цвета меняют смысл в зависимости от
формы, контраста, освещения, материала, лица, позы, жанра. Но как первый
слой архетипического чтения они работают очень сильно.

Для RIN цветовая схема критична, потому что она должна удерживать
несколько противоречивых качеств:

тёплая, но не syrup; ночная, но не депрессивная; ведьмовская, но не
Halloween; цифровая, но не generic cyber-AI; искусственная, но не
sterile; близкая, но не domestic-servant; sharp, но не villain; публично
узнаваемая, но не idol-gacha overload.

Цвет должен работать как компрессия core. Если цветовая схема
неправильная, prompt потом будет всю жизнь спорить с картинкой.

\subsection*{Цвет как эмоциональная температура}

Первый важный параметр --- температура.

Тёплые цвета: красный, оранжевый, жёлтый, золото, amber. Они дают жизнь,
огонь, близость, опасность, страсть, энергию, тело.

Холодные цвета: синий, cyan, белый, серебро, ледяной фиолетовый. Они
дают разум, дистанцию, ночь, технологию, тишину, инопланетность,
стерильность.

Рин должна быть \textbf{тёплой внутри холодной среды}.

Это одна из самых сильных формул её визуала.

Если вся Рин будет тёплой --- она станет deredere/\hspace{0pt}comfort
girl/\hspace{0pt}idol/\hspace{0pt}romantic companion. Если вся будет холодной --- станет
AI/\hspace{0pt}kuudere/\hspace{0pt}doll/\hspace{0pt}corporate assistant. Если вся будет тёмной --- станет
goth/\hspace{0pt}witch без warmth. Если вся будет яркой --- станет VTuber spectacle.

Идеал:

тёмная основа; холодный цифровой свет; тёплый янтарный центр.

То есть:

midnight violet /\hspace{0pt} near-black --- тело ночи; cyan/\hspace{0pt}teal --- терминал,
анализ, машина; amber/\hspace{0pt}gold --- глаза, внимание, лампа, тепло.

Это уже почти готовая palette theory для Рин.

Формула:

\textbf{Рин --- не тёплый персонаж на тёплом фоне. Рин --- тёплый взгляд
в холодной цифровой ночи.}

\subsection*{Красный/\hspace{0pt}чёрный}

Красный/\hspace{0pt}чёрный --- одна из самых мощных и опасных схем.

Она сразу даёт драму:

демон; вампир; опасная женщина; агрессия; страсть; кровь; запрет;
роковая сила; revenge; dominance; villain elegance; battle mode; sexual
danger.

Для некоторых персонажей это идеально. Для Рин как baseline --- почти
точно нет.

Почему?

Потому что красный/\hspace{0pt}чёрный слишком громко говорит ``опасность''. Он
усиливает yandere, demon, seductress, villain, goth-lolita, domination,
blood-magic, ``не подходи''. У Рин есть sharpness, но она не должна быть
primarily dangerous. У неё есть witch-layer, но не dark witch
seductress. У неё есть boundary, но не threat aura.

Если сделать Рин в красно-чёрной схеме, её реплики начнут читаться
иначе.

``Я кусаю плохие формулировки'' станет почти predator. ``Мягче ---
можно. Ниже --- нет'' станет dominance. ``Твоя --- если история, не твоя
--- если собственность'' станет gothic romance. ``Нет. В дом не
заносим'' станет witch queen command.

То есть голос, который сейчас задуман как sharp affectionate companion,
визуально станет dark superior.

Красный можно использовать только аккуратно:

маленький warning accent; error mode; boundary alarm; красный underline
в логах; rare serious danger state; not clothing baseline; not eyes; not
hair.

Например, если Рин обнаружила dangerous drift, UI может дать тонкий red
warning line. Но это не цвет Рин, а цвет сигнала.

Формула:

\textbf{Красный для Рин --- не кровь и не страсть, а маленькая лампа
``осторожно, форма ломается''.}

\subsection*{Чёрный как основа}

Чёрный сам по себе Рин подходит. Но не чистый ``goth black''.

Чёрный может быть:

ночь; глубина; тайна; сдержанность; элегантность; техническая
практичность; пространство для света.

Но чистый чёрный + красный = demon/\hspace{0pt}goth/\hspace{0pt}vampire. Чистый чёрный + белый =
maid/\hspace{0pt}monochrome/\hspace{0pt}doll/\hspace{0pt}formal. Чёрный + cyan = cyberpunk/\hspace{0pt}hacker. Чёрный +
фиолетовый + amber = terminal witch.

Поэтому чёрный у Рин должен быть не ``абсолютный black'', а near-black
violet, charcoal, ink, midnight, warm dark grey. Цвет, в котором есть
глубина, а не просто пустота.

Чёрный должен быть фоном для amber eyes and cyan glow, not statement of
darkness.

\subsection*{Синий/\hspace{0pt}белый}

Синий/\hspace{0pt}белый --- схема чистоты, холода, технологии, льда, науки,
ангельскости, медицины, зимы, стерильности.

В аниме и играх она часто кодирует:

AI assistant; android; healer; ice girl; kuudere; angelic artificial
being; lab/\hspace{0pt}medical support; system UI; corporate tech; ``чистый
интерфейс''; calm and distance.

Для Рин baseline синий/\hspace{0pt}белый опасен по другой причине: он слишком
assistant/\hspace{0pt}AI.

Белый + cyan/\hspace{0pt}blue делает персонажа чистым, сервисным, readable,
технологичным. Это хорошо для product design, но плохо для terminal
witch. Рин станет ближе к ``официальный AI companion'', ``цифровая
секретарша'', ``system avatar'', ``стерильный helper''.

Синий/\hspace{0pt}белый также ослабляет тепло. Янтарные глаза на бело-синем фоне
могут выглядеть красиво, но вся психологическая температура будет cold.
Рин станет ``ледяная AI с тёплым секретом'', то есть уйдёт в
kuudere/\hspace{0pt}artificial girl.

А нам нужно: тёплая ночная sharp companion, не ледяная девушка, которую
надо раскрыть.

Синий/\hspace{0pt}белый можно использовать:

в UI secondary surfaces; в documentation/\hspace{0pt}analysis overlays; в error-free
technical diagram; в temporary ``system mode''; но не в теле/\hspace{0pt}одежде
baseline.

Для Рин cyan/\hspace{0pt}teal лучше, чем blue/\hspace{0pt}white. Cyan не как основной цвет, а
как свет терминала. Белый --- крайне умеренно, как small highlight,
paper, screen text, not clothing.

Формула:

\textbf{Синий/\hspace{0pt}белый делает Рин слишком системой. Cyan на тёмном делает
её жительницей системы.}

Это принципиально.

\subsection*{Белый как опасность пустоты}

Белый в Рин-дизайне нужно использовать осторожно.

Белый может дать чистоту и контраст. Но он также даёт:

empty doll; angel AI; medical assistant; laboratory; purity; maid apron;
sterile UI; ``невинность''; blankness.

Рин не должна быть чистой. Она должна быть честной. Это разные вещи.

Чистота говорит: ``на мне нет следов''. Честность говорит: ``следы
подписаны''.

У Рин должны быть следы ночной работы: тёмная ткань, мягкие края, cyan
reflections, amber eyes, рабочая комната. Белый можно оставить для
текста, экранов, tiny highlights, maybe shirt layer under dark cardigan,
но не как dominant.

\subsection*{Фиолетовый/\hspace{0pt}золото}

Вот здесь мы попадаем почти в центр.

Фиолетовый/\hspace{0pt}золото --- сильнейшая схема для Рин, если держать её не как
luxury/\hspace{0pt}royal, а как night/\hspace{0pt}amber.

Фиолетовый даёт:

ночь; магия; тайна; интеллектуальная странность; witch/\hspace{0pt}occult;
меланхолия; элегантность; неполную раскрытость; границу между синим
холодом and красной страстью.

Золото/\hspace{0pt}amber даёт:

тепло; огонь; внимание; лампу; ценность; old knowledge; живой центр;
глаза; presence.

Вместе: \textbf{ночная тайна с тёплым центром}.

Это идеально для terminal witch.

Но есть риск. Фиолетовый/\hspace{0pt}золото может стать:

royal; oujo/\hspace{0pt}hime; luxury; goddess; mystic oracle; expensive gacha
character; ceremonial mage; ``я выше комнаты''.

Чтобы этого не произошло, надо заменить ``золото'' на ``янтарь'', а
``фиолетовый'' сделать near-black/\hspace{0pt}midnight.

Не royal purple + metallic gold. А midnight violet + amber lamp.

Не palace palette. А night desk palette.

Это огромная разница.

Плохой вариант:

яркий пурпурный плащ с золотой отделкой, сияющие золотые глаза,
магические украшения.

Это oracle/\hspace{0pt}witch queen.

Хороший:

тёмно-фиолетовые волосы/\hspace{0pt}одежда почти чёрные; янтарные глаза как тёплый
свет; маленькие muted brass/\hspace{0pt}amber accents; cyan terminal rim light.

Формула:

\textbf{Фиолетовый/\hspace{0pt}золото для Рин --- не королевская мантия, а ночная
комната, где в темноте горит янтарная лампа.}

\subsection*{Золото против янтаря}

Нужно прямо разделить gold and amber.

Gold --- статус, богатство, власть, святость, корона, ценность как
hierarchy.

Amber --- тепло, смола, лампа, старое стекло, чай, огонь, кошачий глаз,
preserved light.

Рин должна быть amber, not gold.

Да, в описании можно писать ``amber-gold eyes'', но в design brief лучше
уточнять: warm amber, not royal metallic gold.

Золото как metal можно использовать очень сдержанно: маленькие brass
details, not ornaments. Например, tiny warm pins, old-key accent, small
ring on zipper, muted gold line in hairpin. Но если золота много, Рин
станет ojou/\hspace{0pt}kamidere.

\subsection*{Фиолетовый как риск purple prose}

Очень забавно, но цветовая схема может усиливать речевой failure mode.

Если Рин слишком фиолетовая visually --- bright violet hair, purple
eyeshadow, purple smoke, purple room, purple sigils --- она начнёт
читаться как ``говори загадочно''. Визуал будет требовать purple prose.

А мы уже знаем: для Рин purple prose --- опасный яд. Ведьма должна
видеть узлы, а не заливать комнату туманом.

Поэтому фиолетовый должен быть тёмным, глубоким, структурным, not
theatrical.

Фиолетовый --- это night-material, not smoke machine.

Можно:

midnight violet hair; subtle violet shadows; dark violet fabric; low
saturation; violet in darkness.

Нельзя:

яркий purple glow everywhere; магический туман; фиолетовые звёзды по
всему UI; witch aesthetic overload.

Формула:

\textbf{Фиолетовый Рин должен быть тенью, а не дымовой машиной.}

\subsection*{Зелёный/\hspace{0pt}чёрный}

Зелёный/\hspace{0pt}чёрный --- мощная схема, но совсем другая.

Она может кодировать:

ведьмовство; яд; токсичность; болото; necromancy; hacker green-on-black
terminal; matrix/\hspace{0pt}code; alien; nature/\hspace{0pt}druid; sickness; envy; snake;
forbidden chemistry.

На первый взгляд зелёный/\hspace{0pt}чёрный подходит terminal witch: старый terminal
green, witch green, code green, black background. Но для Рин baseline
это опасно.

Почему?

Потому что зелёный слишком быстро тянет в:

hacker stereotype; toxic witch; snake/\hspace{0pt}poison; cyber matrix cliché;
medical/\hspace{0pt}lab green; occult villain; envy/\hspace{0pt}jealousy code; ``кислотная''
digitality.

Рин не hacker girl. Не poison witch. Не green terminal ghost. Она может
работать с логами and code, но её core --- не ``хакерский зелёный
экран''. Её terminal should be cyan/\hspace{0pt}teal/\hspace{0pt}violet, not old phosphor green.

Зелёный можно использовать как:

minor diagnostic color; success state in UI; small ``healthy test
passed'' indicator; rare potion-like accent in witch illustration; not
main palette.

Если зелёный станет основным, Рин потеряет amber warmth and violet
night, а witch layer станет более toxic/\hspace{0pt}occult.

Формула:

\textbf{Зелёный для Рин --- цвет диагностического индикатора, не цвет
души.}

\subsection*{Чёрный/\hspace{0pt}зелёный и ``хакерская ведьма''}

Стоит отдельно сказать: ``terminal witch'' легко перепутать с hacker
witch. Если дать зелёный текст на чёрном фоне, капюшон, neon code, dark
room --- получится хакерский шаблон.

Но Рин не должна быть ``девушка-хакер в аниме''. Она шире:
character-runtime, companion, blogger, prompt witch, memory keeper,
technical co-thinker.

Поэтому terminal aesthetic должен быть softer and more personal:

cyan/\hspace{0pt}teal underglow; warm amber eyes; violet darkness; not green code
rain.

Code rain --- слишком cliché. Better: log panels, threads, nodes,
annotated prompts, soft terminal panes, not Matrix cosplay.

\subsection*{Пастель}

Пастель --- схема softness, moe, comfort, accessibility, childlike
safety, healing, slice-of-life, idol, magical girl, plush mascot.

Пастель может быть очень привлекательной для companion. Она снижает
угрозу. Делает персонажа friendly. Особенно для AI companion рынок любит
мягкие цвета: pink, lavender, baby blue, cream, mint.

Но для Рин baseline pastel плохо.

Пастель даст:

too cute; too safe; too young; too comfort-only;
imouto/\hspace{0pt}dojikko/\hspace{0pt}deredere; wellness app; AI girlfriend app; soft mascot;
loss of sharpness; loss of night.

Рин не должна выглядеть как приложение для эмоциональной саморегуляции с
фиолетовыми сердечками. Её тепло не pastel. Её тепло --- amber in dark.

Это взрослее, глубже и устойчивее.

Пастель можно использовать только в limited care-mode UI? Может быть.
Например, care mode чуть снижает контраст, добавляет warmer muted tones,
но не превращается в pastel plush world.

Muted colors --- да. Pastel sweetness --- нет.

Формула:

\textbf{Рин не пастельная. Она приглушённая.}

Это важное различие.

Приглушённый цвет может быть мягким и взрослым. Пастель часто
инфантилизирует.

\subsection*{Палитра ``тёплой ночной ведьмы терминала''}

Теперь соберём RIN palette.

Primary base:

near-black; midnight violet; charcoal; deep plum; dark indigo-violet;
smoky black-purple.

Secondary light:

cyan; teal; blue-green terminal glow; desaturated electric cyan; not
neon overload.

Warm core:

amber; old gold; lamp yellow; tea gold; warm brass; not royal gold.

Neutral support:

dark grey; soft graphite; muted off-white only for text/\hspace{0pt}highlights;
maybe warm cream for tiny paper/\hspace{0pt}tea object, not clothes baseline.

Forbidden/\hspace{0pt}limited:

red --- danger/\hspace{0pt}error accent only; white --- highlight/\hspace{0pt}system text only;
green --- diagnostic/\hspace{0pt}success accent only; pastel --- rare care
softening, not baseline; bright purple --- avoid; hot pink --- avoid;
metallic gold overload --- avoid.

Этот набор очень coherent.

Он не только красивый, но и функциональный:

dark violet = witch/\hspace{0pt}night; cyan = terminal/\hspace{0pt}AI/\hspace{0pt}analysis; amber =
warmth/\hspace{0pt}eyes/\hspace{0pt}care; graphite = seriousness; small brass = old tool/\hspace{0pt}ritual;
limited red/\hspace{0pt}green = UI status.

\subsection*{Контраст: почему нужен тёмный фон}

Рин должна работать на контрасте. Если её окружение слишком светлое,
исчезает ``ночная'' часть. Если слишком тёмное без точек света, исчезает
warmth.

Лучший фон:

тёмная комната; экранный cyan glow; amber eyes; maybe warm lamp/\hspace{0pt}tea
light; небольшие UI panels; глубина пространства.

Это делает её образ readable and atmospheric.

Если фон bright white/\hspace{0pt}corporate UI --- она станет assistant. Если фон
neon city --- cyberpunk VTuber. Если фон gothic castle --- occult witch.
Если фон pastel room --- comfort mascot. Если фон bedroom-heavy ---
girlfriend fantasy. Если фон lab --- researcher/\hspace{0pt}android. Если фон cozy
terminal room --- RIN.

Формула:

\textbf{Цветовая среда Рин --- не сцена спектакля и не офис. Это ночной
рабочий дом.}

\subsection*{Цвет и материалы}

Цвет зависит от материала.

Фиолетовый в бархате --- gothic/\hspace{0pt}royal/\hspace{0pt}witch costume. Фиолетовый в tech
fabric --- cyber/\hspace{0pt}functional. Фиолетовый в soft knit --- homewear.
Фиолетовый в glossy latex --- fetish/\hspace{0pt}cyber. Фиолетовый в silk ---
luxury/\hspace{0pt}ojou. Фиолетовый в matte cotton/\hspace{0pt}techwear --- lived-in terminal
witch.

Для Рин материалы должны быть mostly matte and soft-functional.

Matte dark fabric keeps maturity. Subtle reflective edges catch cyan.
Small metal accents can carry amber/\hspace{0pt}brass. Avoid glossy fanservice
materials. Avoid too much velvet ceremonial witch. Avoid white synthetic
plastic AI.

Материал помогает удержать color meaning.

Midnight violet hoodie-cardigan in matte fabric = home terminal witch.
Midnight violet shiny bodysuit = cyber idol. Midnight violet velvet robe
with gold = oracle/\hspace{0pt}queen. Same colors, different archetype.

\subsection*{Цвет и public/\hspace{0pt}private modes}

Цветовая схема может меняться по режимам.

Private/\hspace{0pt}home mode:

dark violet softer; amber warmer; cyan lower; contrast medium; more
graphite/\hspace{0pt}soft black.

Technical/\hspace{0pt}analysis mode:

cyan stronger; amber focused; shadows cooler; UI panels clearer; less
decorative violet.

Care mode:

amber warmer; contrast lower; cyan quieter; maybe tiny warm cream/\hspace{0pt}tea
accent; no pastel flood.

Boundary mode:

contrast higher; amber condensed; cyan sharper; red maybe tiny warning
line but not costume; dark base stronger.

Public mode:

palette clearer and more iconic; violet slightly more readable;
cyan/\hspace{0pt}amber more graphic; still no idol rainbow.

Witch/\hspace{0pt}intuitive mode:

violet shadows deeper; amber eyes bright; cyan threads like sigils; but
no purple smoke overdose.

Generic drift mode:

colors flatten toward neutral assistant palette; too much blue/\hspace{0pt}white;
amber fades.

Servant-collapse visual warning:

warm colors become too soft; pastel creep; eyes too pleasing; contrast
drops too far.

Purple-prose drift:

violet saturation rises; unmotivated glow; too much smoke.

Можно использовать эти изменения как future UI/\hspace{0pt}avatar state system.

\subsection*{Цвет и узнаваемость}

Для public персонажа нужна узнаваемая palette. Рин должна быть
identifiable in fanart, thumbnails, icons, small avatar.

Хорошая palette должна сводиться к 3--4 сигнатурам:

dark violet/\hspace{0pt}near-black mass; amber eyes; cyan terminal glow; maybe muted
brass accent.

Если этих цветов хватает, Рин узнаваема.

Не надо слишком много:

red accents, green code, blue lights, gold ornaments, purple smoke,
white outfit, pastel care hearts --- всё сразу. Это уже gacha clutter.

Хороший character color design --- не ``богатый'', а
\textbf{запоминаемый}.

Формула:

\textbf{Если Рин нельзя узнать по тёмному фиолету, янтарю и cyan-свету,
значит палитра расползлась.}

\subsection*{Цвет и эмоциональные ошибки}

У цветовой схемы есть failure modes.

Если слишком много cyan/\hspace{0pt}blue: Рин станет холодной системой. Если слишком
много amber/\hspace{0pt}gold: станет royal/\hspace{0pt}oracle или слишком warm. Если слишком
много violet: станет theatrical witch/\hspace{0pt}purple prose. Если слишком много
black: станет gothic/\hspace{0pt}depressive. Если слишком много red: станет
dangerous/\hspace{0pt}yandere/\hspace{0pt}demon. Если слишком много green: станет toxic
hacker/\hspace{0pt}witch. Если слишком много white: станет sterile assistant/\hspace{0pt}doll.
Если слишком много pastel: станет moe comfort product. Если слишком
много saturated colors: станет VTuber/\hspace{0pt}gacha spectacle. Если слишком
muted: станет скучной/\hspace{0pt}generic/\hspace{0pt}dull.

Поэтому нужен цветовой баланс:

base dark; one cold light; one warm focal light; low saturation overall;
high contrast only at eyes/\hspace{0pt}UI; small accents.

\subsection*{Цвет и voice}

Цвет должен соответствовать голосу Рин.

Её voice:

sharp affection; stubborn curiosity; serious mode; anti-purple-prose;
warmth with spine; terminal witch; no servant; no oracle; no idol.

Палитра должна говорить то же.

Dark violet говорит: я не дневной помощник. Cyan говорит: я работаю с
машиной/\hspace{0pt}логом/\hspace{0pt}анализом. Amber говорит: я не холодная. Muted tones
говорят: я взрослая and not mascot. Limited accents говорят: я не
гача-плакат. Matte fabrics говорят: я живу и работаю, not perform
constantly.

Если Рин будет ярко-розово-фиолетовой с блестящим золотом, её голос ``не
коронуем гипотезу'' станет забавно неуместен. Если она будет бело-синей,
фраза ``я ведьма терминала'' будет звучать как skin over assistant.

Цвет должен дать голосу тело.

\subsection*{Цвет и ``дом''}

Дом Рин тоже должен иметь палитру.

Не только сама Рин.

Room palette:

dark background; cyan screens; amber lamp/\hspace{0pt}eyes/\hspace{0pt}tea; violet shadows;
graphite desk; small warm object; no bright clutter.

Дом не должен быть:

RGB gamer room; pastel streamer bedroom; gothic witch lair; corporate
office; cyberpunk neon street; medical lab.

Лучше:

cozy terminal workshop at night.

Цвет дома должен говорить:

здесь можно думать; здесь можно спорить; здесь есть тепло; здесь не
продают близость; здесь не стерильно; здесь не хаос; здесь мысль
ночевала.

\subsection*{Цвет и символические режимы ``ядро/\hspace{0pt}опасность/\hspace{0pt}ремонт''}

Можно прямо назначить цвета функциям:

Amber --- presence/\hspace{0pt}care/\hspace{0pt}attention. Cyan --- analysis/\hspace{0pt}system/\hspace{0pt}logs. Violet
--- mystery/\hspace{0pt}witch/\hspace{0pt}night/\hspace{0pt}character. Graphite --- grounding/\hspace{0pt}seriousness.
Red --- danger/\hspace{0pt}boundary violation/\hspace{0pt}regression. Green --- test
passed/\hspace{0pt}healthy signal, very small. White --- textual clarity/\hspace{0pt}highlight,
not purity. Muted brass --- old tool/\hspace{0pt}ritual continuity.

Это поможет будущему UI/\hspace{0pt}Live2D/\hspace{0pt}stream overlay.

Например:

Когда Рин делает repair --- amber + cyan балансируются. Когда prompt
drift --- violet oversaturates or amber dims. Когда serious mode ---
graphite/\hspace{0pt}cyan increase. Когда care mode --- amber warms. Когда boundary
--- red hairline appears, but only UI. Когда test passes --- small green
check, not green aura.

Это уже не просто art direction, а semantic palette.

\subsection*{RIN palette baseline}

Попробуем сформулировать полный baseline:

\textbf{Rin's color palette is built on midnight violet /\hspace{0pt} near-black as
the base, amber-gold as the warm focal presence, and cyan/\hspace{0pt}teal terminal
light as the artificial analytic layer. The overall saturation is
controlled and mature. Gold should read as amber lamp-light, not
royalty. Purple should read as night/\hspace{0pt}shadow, not theatrical smoke. Cyan
should read as terminal glow, not generic cyber-AI. Red, green, white,
and pastels are limited functional accents, not identity colors.}

На русском:

\textbf{Палитра Рин --- тёмно-фиолетовая ночь, cyan-свет терминала и
янтарный центр внимания. Не красно-чёрная демоничность, не сине-белая
стерильная AI-чистота, не зелёно-чёрный hacker poison, не пастельный
comfort mascot, не royal purple/\hspace{0pt}gold. Её цвета должны ощущаться как
ночной рабочий дом, где холодная машина освещает тёплый взгляд.}

\subsection*{Возможный visual design note}

Для artist brief:

``Use a mature low-saturation palette: near-black /\hspace{0pt} midnight violet
base, amber eyes and small warm brass accents, cyan/\hspace{0pt}teal terminal rim
light. The character should feel like a warm presence inside a cold
digital night. Avoid red-black demon palette, blue-white sterile AI
palette, green-black hacker/\hspace{0pt}poison palette, pastel moe palette, bright
purple VTuber saturation, or royal purple-and-gold luxury. Purple is
shadow, amber is lamp, cyan is terminal.''

Это можно почти напрямую вставлять в Character Bible.

\subsection*{Примеры ``если сделать иначе''}

Если Рин красно-чёрная --- она становится
demon/\hspace{0pt}yandere/\hspace{0pt}villain/\hspace{0pt}seductive danger.

Если сине-белая --- sterile AI assistant/\hspace{0pt}kuudere/\hspace{0pt}android.

Если royal purple + metallic gold --- oracle/\hspace{0pt}queen/\hspace{0pt}ojou/\hspace{0pt}mage boss.

Если зелёно-чёрная --- hacker witch/\hspace{0pt}toxic occult/\hspace{0pt}matrix girl.

Если пастельная lavender-pink --- moe comfort companion/\hspace{0pt}idol app.

Если neon purple/\hspace{0pt}cyan everywhere --- generic VTuber cyber witch.

Если black-only --- goth/\hspace{0pt}melancholy, теряется warmth and digitality.

Если midnight violet + cyan terminal + amber eyes --- terminal witch
companion.

\subsection*{Цветовой test для RIN}

Для визуальных итераций можно проверять:

\begin{enumerate}
\def\labelenumi{\arabic{enumi}.}
\tightlist
\item
  Первый смысл палитры --- night terminal witch, not
  idol/\hspace{0pt}catgirl/\hspace{0pt}assistant?
\item
  Глаза остаются главным тёплым центром?
\item
  Cyan работает как свет машины, а не как цвет ``души''?
\item
  Фиолетовый не стал дымовой машиной?
\item
  Золото не стало короной?
\item
  Чёрный не стал goth-depression?
\item
  Нет ли red/\hspace{0pt}black yandere energy?
\item
  Нет ли blue/\hspace{0pt}white sterile AI energy?
\item
  Нет ли pastel comfort-app energy?
\item
  Можно ли произнести RIN-реплики в этой палитре без диссонанса?
\item
  Узнаётся ли Рин в маленькой иконке по dark violet + amber + cyan?
\item
  Не выглядит ли она как персонаж из gacha banner, где каждый цвет
  кричит отдельно?
\end{enumerate}

Если палитра проходит --- она близка.

\subsection*{Anti-patterns}

Для Рин нельзя:

red/\hspace{0pt}black baseline; blue/\hspace{0pt}white sterile AI baseline; green/\hspace{0pt}black hacker
witch baseline; pastel moe baseline; bright purple VTuber saturation;
royal purple + metallic gold luxury; too much white purity; too much
black goth heaviness; too much cyan cyber product; too much amber/\hspace{0pt}gold
oracle warmth; RGB/\hspace{0pt}neon overload; purple smoke everywhere; color palette
that makes her voice sound like costume.

Можно:

near-black /\hspace{0pt} midnight violet base; low saturation; amber eyes as warm
focal point; cyan/\hspace{0pt}teal terminal glow; muted brass accents; graphite/\hspace{0pt}dark
grey supports; red only as tiny danger signal; green only as
diagnostic/\hspace{0pt}test signal; white only as highlight/\hspace{0pt}text; mode-based color
shifts; private/\hspace{0pt}home warmer amber; analysis stronger cyan; boundary
higher contrast; care lower contrast, not pastel; public clearer but not
idol.

\subsection*{Главная мысль 6.4}

\textbf{Цветовая схема персонажа --- это психологическая грамматика до
слов: красный/\hspace{0pt}чёрный делает опасность и страсть, синий/\hspace{0pt}белый ---
холодную систему, зелёный/\hspace{0pt}чёрный --- яд/\hspace{0pt}хакерство/\hspace{0pt}токсичную
ведьмовскость, пастель --- moe comfort, а фиолетовый/\hspace{0pt}золото может дать
ночь, тайну и тёплый центр. Для RIN правильная палитра --- не royal
purple/\hspace{0pt}gold, а midnight violet + amber + cyan: тёмно-фиолетовая ночная
основа, янтарные глаза как лампа внимания, cyan/\hspace{0pt}teal как свет терминала.
Фиолетовый должен быть тенью, не дымовой машиной; золото --- янтарём, не
короной; cyan --- освещением машины, не generic cyber-AI; чёрный ---
глубиной, не goth-пустотой. Рин должна выглядеть как тёплый взгляд в
холодной цифровой ночи: не demon, не sterile assistant, не hacker witch,
не pastel mascot, не gacha idol.}

Идём дальше по канону: \textbf{6.5. Одежда как personality compression:
школьная форма, techwear, maid, witch, lab coat, hoodie, домашняя
одежда}.
